\begin{table}[H]
\centering
\caption{Resultados sobre el conjunto de prueba.}
\label{tab:resultados}
\begin{tabular}{lrrrr}
\toprule
Modelo & Accuracy & Precision macro & Recall macro & F1 macro\\
\midrule
A & 0.8802 & 0.8798 & 0.8814 & 0.8799 \\
B & 0.8784 & 0.8785 & 0.8777 & 0.8775 \\
C & 0.8675 & 0.8713 & 0.8644 & 0.8652 \\
\bottomrule
\end{tabular}
\end{table}

\paragraph{Selección del mejor modelo.}
El Modelo A obtuvo el mayor F1 macro (0.8799), con Accuracy 0.8802. Por tanto, bajo la regla definida antes del experimento, es el modelo seleccionado. La interpretación por clase y los posibles indicios de sobreajuste deben contrastarse con las matrices de confusión y las curvas presentadas.

\paragraph{Análisis por clase.} En el modelo ganador, la clase con mayor F1 fue sunflowers (0.9048) y la de menor F1 fue roses (0.8317). La brecha absoluta entre F1 macro de validación y prueba fue 0.0219; este valor, junto con las curvas de aprendizaje, permite valorar la generalización sin basarse solamente en Accuracy.
